\chapter{}
\label{chap:2}

\section{Умова}
Нехай $I$ -- ідеал кiльця $B_{n}$. Обчислити потужність $\left| I\right|$, якщо:
\begin{enumerate}[label=\alph*)]
    \item $n$ є парним числом та $I = \langle x_{1}x_{2}, \, x_{3}x_{4}, \, \ldots, \, x_{n}x_{n-1} \rangle$;
    \item $n \geq 2$ та $I = \langle x_{1}x_{2}, \, x_{1}x_{2}x_{3}, \, \ldots , \, x_{1}x_{2} \ldots x_{n}\rangle$.
\end{enumerate}

\section{Розв'язання}

Булеве кільце $B_{n}$ від $n$ змінних, еквівалентно фактор кільцю $\mathbb{F}_2[x_1, \ldots, x_n] / \langle x_1^2 \oplus x_1, \ldots, x_n^2 \oplus x_n \rangle$. 
Тобто кожен його елемент можна представити як многочлен, де кожна зі змінних $x_{i}$ входить не більше чим в першому 
степені, тобто мономи. $\left|B_{n}\right| = 2^{2^{n}}$.

В ідеалі $I \subseteq B_n$ множина нулів складається з усіх точок, в яких кожна функція з ідеалу обертається в нуль:
\begin{equation*}
    V(I) = \left\{a \in \left\{0,1\right\}^n : f(a) = 0, \, \forall \, f \in I\right\}
\end{equation*}

Еквівалентне означення: Для ідеалу $I = \langle g_1, \ldots, g_m \rangle$ достатньо перевірити лише генератори:
\begin{equation*}
    V(I) = \left\{a \in \left\{0,1\right\}^n : g_1(a) = g_2(a) = \ldots = g_m(a) = 0\right\}
\end{equation*}

\noindent Функція $f \in I \Leftrightarrow f(a) = 0$ для всіх $a \in V(I)$.

Отже:
\begin{itemize}
    \item На точках $a$ з $V(I)$ значення функцій з $I$ фіксоване (дорівнює 0)
    \item На решті $2^n - |V(I)|$ точок значення є довільним (0 або 1)
\end{itemize}

Тому потужність ідеалу:
\begin{equation}
    \label{eq:ideal_power}
    |I| = 2^{2^n - |V(I)|}
\end{equation}

\subsection{Пункт а)}

\textbf{Дано:} $n$ парне, $I = \langle x_{1}x_{2}, \, x_{3}x_{4}, \, \ldots, \, x_{n}x_{n-1} \rangle$.

Знайдемо $V(I)$ як перетин множин нулів:
\begin{equation*}
    V(I) = V(x_1 x_2) \cap V(x_3 x_4) \cap \ldots \cap V(x_{n-1} x_n)
\end{equation*}

Для деякої пари $(x_{2k-1}, x_{2k})$:
\begin{equation*}
    V(x_{2k-1} x_{2k}) \overset{def}{=} \left\{a \in \left\{0,1\right\}^n : a_{2k-1} \cdot a_{2k} = 0\right\}
\end{equation*}

З цього можна зробити висновок, що $a_{2k-1} = 0$ або $a_{2k} = 0$. Отже пара $(a_{2k-1}, a_{2k})$ може набувати 
3 значення з 4: $(0, 0)$, $(0, 1)$, $(1, 0)$. Оскільки пари є незалежними між собою, а їх кількість $\frac{n}{2}$, 
кожна має 3 допустимими варіантами, то:
\begin{equation*}
    |V(I)| = 3^{n/2}
\end{equation*}

Тоді для обчислення $|I|$ використаємо вище згадану формулу~\eqref{eq:ideal_power}:
\begin{equation*}
    |I| = 2^{2^n - |V(I)|} = 2^{2^n - 3^{n/2}}
\end{equation*}

\subsection{Пункт б)}

\textbf{Дано:} $n \geq 2$, $I = \langle x_{1}x_{2}, \, x_{1}x_{2}x_{3}, \, \ldots , \, x_{1}x_{2} \ldots x_{n}\rangle$.

Можна сказати, що всі генератори є кратними до першого ($x_1 x_2$):
\begin{equation*}
    x_1 x_2 x_3 = x_3 \cdot (x_1 x_2) \in \langle x_1 x_2 \rangle
\end{equation*}
\begin{equation*}
    x_1 x_2 x_3 x_4 = x_3 x_4 \cdot (x_1 x_2) \in \langle x_1 x_2 \rangle
\end{equation*}
\begin{equation*}
    \vdots
\end{equation*}
\begin{equation*}
    x_1 x_2 \ldots x_n = x_3 x_4 \ldots x_n \cdot (x_1 x_2) \in \langle x_1 x_2 \rangle
\end{equation*}

$\Rightarrow$ (?) всі генератори належать ідеалу $\langle x_1 x_2 \rangle$ (щось я в коректності цього твердження 
сумніваюся):
\begin{equation*}
    I = \langle x_1 x_2 \rangle
\end{equation*}
Тоді $V(I)$ -- множина всіх точок, де $a_1 = 0$ або $a_2 = 0$, оскільки:
\begin{equation*}
    V(I) = V(x_1 x_2) = \left\{a \in \left\{0,1\right\}^n : a_1 \cdot a_2 = 0\right\}
\end{equation*}

Обчислимо $|V(I)|$, але, певно, доведеться формулу включення-виключення застосовувати обчислюючи доповнення спершу:
\begin{equation*}
    \overline{V(I)} = \left\{a : a_1 = 1 \text{ і } a_2 = 1\right\}
\end{equation*}
Це будуть такі точки $a$, де перші дві координати фіксовані значенням $1$, а решта $n - 2$ координат довільні, 
тобто $(1, 1, 0/1, \ldots, 0/1)$. Тоді:
\begin{equation*}
    |\overline{V(I)}| = 2^{n-2}
\end{equation*}
А тоді знайдемо $|V(I)|$, віднявши $|\overline{V(I)}|$ від загальної кількості:
\begin{equation*}
    |V(I)| = 2^n - 2^{n-2} = 2^{n-2}(4 - 1) = 3 \cdot 2^{n-2}
\end{equation*}
Ну і власне застосуємо вище згадану формулу~\eqref{eq:ideal_power}:
\begin{equation*}
    |I| = 2^{2^n - |V(I)|} = 2^{2^n - 3 \cdot 2^{n-2}} = 2^{2^{n-2}}
\end{equation*}
Отже:
\begin{equation*}
    |I| = 2^{2^{n-2}}
\end{equation*}